\subsection{自同构的模}

设 G 是局部紧群，\(\varphi\) 是 G 的自同构，\(\mu\) 是 G 上的左哈尔测度。显然，\(\varphi^{-1}(\mu)\) 也是 G 上的左哈尔测度。因此存在（No.~2, Th. 1）唯一的数 a \textgreater{} 0，使得 \(\varphi^{-1}(\mu) = a\mu\)。根据 No.~2, Th. 1，该数与 \(\mu\) 的选择无关。注意，若一开始取右哈尔测度，例如 \(\Delta_{G}^{-1} \cdot \mu\)（No.~3），也会得到同一个标量 a：因为 \(\varphi^{-1}\) 保持 \(\Delta_{G}\) 不变（No.~3, 评注），有 \(\varphi^{-1}(\Delta_{G}^{-1} \cdot \mu) = \Delta_{G}^{-1} \cdot \varphi^{-1}(\mu) = a\Delta_{G}^{-1} \cdot \mu\)。

定义 4. --- 数 a \textgreater{} 0 满足 \(\varphi^{-1}(\mu) = a\mu\)，称为自同构 \(\varphi\) 的模，记为 mod G \(\varphi\)，或简称为 mod \(\varphi\)。

若 \(f\) 是在 \(\mu\) 下于 \(G\) 可积的函数，则

\[
\int f \big (\varphi^ {- 1} (x) \big) d \mu (x) = (\mathrm{mod} \varphi) \int f (x) d \mu (x).\tag{31}
\]

若 A 是 G 的 \(\mu\)-可积子集，则

\[
\mu (\varphi (\mathrm{A})) = (\mathrm{mod} \varphi) \mu (\mathrm{A}).\tag{32}
\]

特别地，对于 \(s \in \mathbf{G}\)，令 \(i_s\) 为内自同构 \(x \mapsto s^{-1}xs\)。则 \(i_s^{-1} = \delta(s)\gamma(s)\)，因而

\[
i _ {s} ^ {- 1} (\mu) = \delta (s) \mu = \Delta_ {\mathrm{G}} (s) \mu ,
\]

所以

\[
\operatorname{mod} i _ {s} = \Delta_ {\mathrm{G}} (s).\tag{33}
\]

若 G 是离散群或紧群，则其归一化哈尔测度在 G 的每个自同构 \(\varphi\) 下都变为自身，这由结构传递立即可见。因此，离散群或紧群的自同构的模为 1。

命题 4. --- 设 G 是局部紧群，\(\Gamma\) 是拓扑群，且 \(\gamma \mapsto u_{\gamma}\) 是从 \(\Gamma\) 到 G 的自同构群 \(\mathcal{G}\) 的同态，并满足 \((\gamma, x) \mapsto u_{\gamma}(x)\) 是从 \(\Gamma \times G\) 到 G 的连续映射。则映射 \(\gamma \mapsto \operatorname{mod}(u_{\gamma})\) 是 \(\Gamma\) 在 \(\mathbf{R}_{+}^{*}\) 中的连续表示。

这个映射显然是 \(\Gamma\) 在 \(\mathbf{R}_+^*\) 中的（代数）表示；只需证明其连续性。取 \(f \in \mathcal{K}(\mathbf{G})\)，令 \(\mathbf{S}\) 为其支集。

取 \(\gamma_0 \in \Gamma\)，令 \(U\) 为 \(u_{\gamma_0}^{-1}(S)\) 的相对紧邻域。映射 \(\gamma \mapsto u_\gamma\) 是从 \(\Gamma\) 到赋予紧收敛拓扑的 \(\mathcal{G}\) 的连续映射（GT, X, §3, No.~4, Th. 3）；因此，\(u_{\gamma}^{-1}(S) \subset U\) 在 \(\gamma\) 充分接近 \(\gamma_0\) 时成立。No.~1 的引理 1 于是证明 \(\int f(u_\gamma(x)) d\mu(x)\)（其中 \(\mu\) 表示 \(G\) 上的左哈尔测度）连续依赖于 \(\gamma\)；命题得证。

